\documentclass[a4paper,11pt]{article}

\usepackage[utf8]{inputenc}
\usepackage[T1]{fontenc}
\usepackage[provide=*,italian]{babel}
\usepackage[a4paper,margin=2.6cm]{geometry}
\usepackage{lmodern}
\usepackage{microtype}
\usepackage{array}
\usepackage{booktabs}
\usepackage{longtable}
\usepackage{tabularx}
\usepackage{hyperref}
\usepackage{xcolor}
\usepackage{enumitem}
\usepackage{listings}
\usepackage{float}
\usepackage{caption}

\definecolor{titleblue}{HTML}{173F5F}
\definecolor{rulegray}{HTML}{C9CED6}
\definecolor{codebg}{HTML}{F7F8FA}
\definecolor{codeborder}{HTML}{D8DDE6}

\hypersetup{
  colorlinks=true,
  linkcolor=titleblue,
  urlcolor=titleblue,
  pdftitle={Testing e valutazione dell'usabilita - DietiEstates25},
  pdfauthor={Gruppo DietiEstates25}
}

\setlist[itemize]{leftmargin=1.1cm}
\setlist[enumerate]{leftmargin=1.1cm}
\renewcommand{\arraystretch}{1.18}
\captionsetup{font=small,labelfont=bf}

\lstdefinestyle{projectcode}{
  language=Java,
  basicstyle=\ttfamily\small,
  backgroundcolor=\color{codebg},
  frame=single,
  rulecolor=\color{codeborder},
  breaklines=true,
  columns=fullflexible,
  keepspaces=true,
  showstringspaces=false,
  commentstyle=\color{gray},
  keywordstyle=\color{titleblue},
  stringstyle=\color{teal}
}

\begin{document}

\begin{titlepage}
  \centering
  \vspace*{2.5cm}
  {\Large Ingegneria del Software\par}
  \vspace{0.7cm}
  {\large Progetto DietiEstates25\par}
  \vspace{1.8cm}
  {\Huge\bfseries Testing e Valutazione dell'Usabilita\par}
  \vspace{1.2cm}
  {\Large Sezione di documentazione per la consegna di progetto\par}
  \vfill
  \begin{tabular}{rl}
    Gruppo: & DietiEstates25 \\
    Artefatti analizzati: & suite Jest backend, documentazione di progetto, UC2, UC4 \\
    Versione documento: & 30 marzo 2026
  \end{tabular}
  \vfill
  \rule{0.82\textwidth}{0.4pt}\par
  \vspace{0.3cm}
  {\small Il layout del presente documento e stato reso coerente con i PDF gia presenti nella cartella Documentazione: frontespizio dedicato, indice, sezioni numerate e tabelle sintetiche per favorire una consegna omogenea.\par}
\end{titlepage}

\pagenumbering{Roman}
\tableofcontents
\clearpage
\pagenumbering{arabic}

\section{Obiettivo e fonti}
Questo documento raccoglie la documentazione necessaria per il punto \emph{Testing e valutazione dell'usabilita} del progetto DietiEstates25. Il contenuto e stato costruito a partire dagli artefatti realmente presenti nel repository e dai risultati prodotti dalla suite di test backend.

L'obiettivo e duplice:
\begin{itemize}
  \item documentare il codice xUnit per quattro metodi non banali, descrivendo la strategia di progettazione dei test con classi di equivalenza e criteri di copertura;
  \item documentare una valutazione di usabilita completa, composta da expert review sistematica e da un esperimento con utenti reali o potenziali corredato da survey e discussione dei risultati.
\end{itemize}

Le fonti utilizzate sono le seguenti:
\begin{itemize}
  \item \texttt{dietiestates-backend/test/backend.test.ts};
  \item \texttt{dietiestates-backend/jest.config.ts};
  \item \texttt{dietiestates-backend/package.json};
  \item \texttt{Progettazione/Test-Doc.md};
  \item \texttt{Progettazione/Requisiti.md};
  \item \texttt{Progettazione/UC2 Caricamento Immobili.md};
  \item \texttt{Progettazione/UC4 Gestione Offerte su.md};
  \item i PDF gia presenti nella cartella \texttt{Documentazione}, usati come riferimento per struttura e tono editoriale.
\end{itemize}

\section{Testing xUnit}

\subsection{Tecnologia e impostazione della suite}
La suite di unit test del backend e implementata con \textbf{Jest} in ambiente Node.js, tramite preset \texttt{ts-jest}. Dall'analisi della configurazione emergono le seguenti scelte:
\begin{itemize}
  \item esecuzione in ambiente \texttt{node};
  \item localizzazione della suite nella cartella \texttt{test/};
  \item raccolta della coverage sui file \texttt{src/**/*.ts};
  \item esclusione di \texttt{src/server.ts} e dell'intera cartella \texttt{src/config/} dal calcolo della copertura.
\end{itemize}

Lo script effettivamente usato per l'esecuzione e il seguente:

\begin{lstlisting}[style=projectcode]
npm test
// equivalente a: jest --runInBand --detectOpenHandles
\end{lstlisting}

La scelta di Jest e coerente con un'impostazione xUnit: ogni metodo viene testato in modo isolato, con fixture leggere, mocking delle dipendenze esterne e asserzioni mirate sul comportamento osservabile.

\subsection{Metodi selezionati}
La suite attuale copre quattro metodi che rispettano il requisito di essere non banali e di ricevere almeno due parametri formali. La Tabella~\ref{tab:metodi} ne riassume la natura.

\begin{table}[H]
\centering
\begin{tabularx}{\textwidth}{@{}c l l X@{}}
\toprule
\# & Metodo & Modulo & Motivazione della scelta \\
\midrule
1 & \texttt{updateStatoOfferta(idOfferta, nuovoStato)} & \texttt{OffertaDAO} & Aggiornamento persistente + mapping DTO con conversioni di tipo \\
2 & \texttt{assignUserToAgency(idUtente, idAgenzia)} & \texttt{UtenteDAO} & UPDATE con WHERE condizionato su ruolo e gestione null \\
3 & \texttt{createOfferta(idImmobile, idUtente, prezzoOfferto)} & \texttt{OffertaDAO} & Inserimento con valori di default di dominio e normalizzazione del record DB \\
4 & \texttt{getNearbyPlaces(lat, lon, radius?)} & \texttt{utils/geoapify} & Chiamate HTTP multiple, lettura variabili d'ambiente e aggregazione risultati \\
\bottomrule
\end{tabularx}
\caption{Metodi non banali coperti dai test xUnit backend.}
\label{tab:metodi}
\end{table}

\subsection{Strategia di progettazione dei test}
La progettazione dei test adotta un approccio \textbf{white box}. Tale scelta e giustificata dal fatto che i metodi verificati sono concentrati nel backend e dipendono da dettagli implementativi noti: query SQL costruite dinamicamente, gestione di campi opzionali, mapping di record provenienti dal database e integrazione con servizi esterni.

Le tecniche principali impiegate sono:
\begin{itemize}
  \item \textbf{classi di equivalenza}, per partizionare il dominio degli input in insiemi significativi di casi validi e non validi;
  \item \textbf{statement coverage}, per verificare il percorso nominale di esecuzione e il corretto completamento delle operazioni previste;
  \item \textbf{branch coverage}, nei metodi con diramazioni esplicite o condizioni di errore;
  \item \textbf{boundary analysis}, dove sono presenti parametri opzionali o valori di default rilevanti;
  \item \textbf{mocking delle dipendenze esterne}, tramite mock di PostgreSQL e Axios per eliminare dipendenze da rete e database reale.
\end{itemize}

Questa impostazione consente test deterministici, veloci e ripetibili, caratteristica essenziale in una pipeline di verifica continua.

\subsection{Analisi dei quattro metodi testati}

\subsubsection{\texttt{OffertaDAO.updateStatoOfferta(idOfferta, nuovoStato)}}
Il metodo aggiorna lo stato di un'offerta nel database e converte la riga restituita nel corrispondente DTO applicativo. La logica non e banale perche unisce persistenza, parametrizzazione SQL e normalizzazione di campi numerici provenienti dal DB come stringhe.

\paragraph{Classi di equivalenza.}
\begin{itemize}
  \item CE1: offerta esistente e nuovo stato valido, ad esempio \texttt{(42, 'Accettata')};
  \item CE2: offerta inesistente, con query che non restituisce righe;
  \item CE3: offerta esistente con campi DB da convertire, ad esempio \texttt{prezzoofferto = '200000'}.
\end{itemize}

\paragraph{Copertura adottata.}
Per questo metodo si applica soprattutto \textbf{statement coverage} sul percorso nominale, con verifica esplicita della query emessa e del mapping del risultato.

\paragraph{Evidenza dal test.}
\begin{lstlisting}[style=projectcode]
const updated = await updateStatoOfferta(42, 'Accettata');

expect(pool.query).toHaveBeenCalledWith(
  'UPDATE offerta SET stato = $1 WHERE idofferta = $2 RETURNING *',
  ['Accettata', 42]
);
expect(updated.prezzoOfferto).toBe(200000);
\end{lstlisting}

\subsubsection{\texttt{UtenteDAO.assignUserToAgency(idUtente, idAgenzia)}}
Il metodo assegna un utente (Agente o Supporto) a un'agenzia, eseguendo una \texttt{UPDATE} con clausola \texttt{WHERE} condizionata sul ruolo. La non banalita risiede nella verifica del ruolo e nella possibilita che l'utente non sia assegnabile (caso in cui il metodo restituisce \texttt{null}).

\paragraph{Classi di equivalenza.}
\begin{itemize}
  \item CE1: utente con ruolo Agente o Supporto, assegnazione riuscita;
  \item CE2: utente inesistente o con ruolo non consentito, restituisce \texttt{null}.
\end{itemize}

\paragraph{Copertura adottata.}
Qui e rilevante la \textbf{branch coverage}: il test copre sia il ramo nominale (assegnazione avvenuta), sia il ramo di mancata assegnazione (\texttt{null}). Viene inoltre verificata la correttezza della query e dei parametri passati.

\paragraph{Evidenza dal test.}
\begin{lstlisting}[style=projectcode]
const updated = await assignUserToAgency(7, 3);
expect(updated?.idUtente).toBe(7);
expect(updated?.idAgenzia).toBe(3);

const result = await assignUserToAgency(999, 3);
expect(result).toBeNull();
\end{lstlisting}

\subsubsection{\texttt{OffertaDAO.createOfferta(idImmobile, idUtente, prezzoOfferto)}}
Il metodo incapsula la creazione di una nuova offerta, valorizzando automaticamente parametri di dominio come \texttt{offertaManuale = false} e \texttt{idOffertaOriginale = null}. Anche in questo caso e presente mapping dal formato DB al DTO applicativo.

\paragraph{Classi di equivalenza.}
\begin{itemize}
  \item CE1: creazione standard con input validi e positivi;
  \item CE2: offerta originale assente, che deve risultare \texttt{null};
  \item CE3: valori numerici restituiti dal DB come stringhe, da convertire.
\end{itemize}

\paragraph{Copertura adottata.}
Si applica \textbf{statement coverage} sul percorso nominale e controllo esplicito del payload inviato all'\texttt{INSERT}.

\paragraph{Evidenza dal test.}
\begin{lstlisting}[style=projectcode]
const offerta = await createOfferta(50, 20, 90000);

expect(pool.query).toHaveBeenCalledWith(
  expect.stringContaining('INSERT INTO offerta'),
  [50, 20, 90000, false, null]
);
expect(offerta.offertaManuale).toBe(false);
\end{lstlisting}

\subsubsection{\texttt{getNearbyPlaces(lat, lon, radius?)}}
Il metodo interroga il servizio esterno Geoapify per piu categorie di punti di interesse e aggrega i risultati. La funzione e non banale perche legge una variabile d'ambiente, invoca tre richieste HTTP distinte e deve gestire il caso di chiave assente.

\paragraph{Classi di equivalenza.}
\begin{itemize}
  \item CE1: chiave API presente e risposte valide dal servizio esterno;
  \item CE2: chiave API assente, con eccezione immediata;
  \item CE3: una o piu categorie senza risultati;
  \item CE4: raggio passato esplicitamente e diverso dal default.
\end{itemize}

\paragraph{Copertura adottata.}
Il metodo richiede \textbf{branch coverage} per coprire il ramo di errore e il flusso nominale. E inoltre presente una forma di \textbf{boundary analysis} sul parametro opzionale \texttt{radius}.

\paragraph{Evidenza dal test.}
\begin{lstlisting}[style=projectcode]
const places = await getNearbyPlaces(45, 9, 500);

expect(axios.get).toHaveBeenCalledTimes(3);
expect(places).toHaveLength(2);
await expect(getNearbyPlaces(45, 9)).rejects.toThrow('Chiave Geoapify mancante');
\end{lstlisting}

\subsection{Riepilogo della copertura progettata}
\begin{table}[H]
\centering
\begin{tabularx}{\textwidth}{@{}l c c c c@{}}
\toprule
Metodo & Equivalenza & Statement & Branch & Note \\
\midrule
\texttt{updateStatoOfferta} & Si & Si & Parziale & Verifica query e mapping DTO \\
\texttt{assignUserToAgency} & Si & Si & Si & Coperto anche il caso di mancata assegnazione \\
\texttt{createOfferta} & Si & Si & Parziale & Verifica valori di default di dominio \\
\texttt{getNearbyPlaces} & Si & Si & Si & Coperti chiave assente e raggio esplicito \\
\bottomrule
\end{tabularx}
\caption{Strategie di copertura adottate nella suite xUnit.}
\label{tab:coverage-strategy}
\end{table}

\subsection{Allineamento completo con Test-Doc}
Questa sezione riporta in forma esplicita tutti gli elementi presenti in \texttt{Progettazione/Test-Doc.md}, in modo che il presente file \texttt{.tex} contenga integralmente il contenuto richiesto.

\subsubsection{Tipologia di test: White box (dettaglio operativo)}
Tutti i test adottati sono \textbf{white box}. In particolare:
\begin{itemize}
  \item e noto il testo delle query SQL prodotte (verifica di UPDATE parametrizzate);
  \item e noto il comportamento di \texttt{insertOfferta} sui parametri \texttt{offertaManuale} e \texttt{offertaOriginaleId -> null};
  \item e noto il flusso di \texttt{getNearbyPlaces} (3 categorie fisse e lettura di \texttt{process.env.GEOAPIFY\_KEY} a runtime);
  \item sono noti i branch interni di \texttt{assignUserToAgency} (check su \texttt{result.rows[0]} e clausola \texttt{Ruolo IN (...)}).
\end{itemize}

Il mocking di \texttt{pool} (PostgreSQL) e \texttt{axios} isola ogni unita dalle dipendenze esterne, rendendo i test deterministici e ripetibili.

\subsubsection{Metodi testati (firma e tipo)}
\begin{table}[H]
\centering
\small
\begin{tabular}{@{}p{0.04\textwidth}p{0.20\textwidth}p{0.15\textwidth}p{0.43\textwidth}p{0.10\textwidth}@{}}
\toprule
\# & Metodo & Modulo & Firma & Tipo \\
\midrule
1 & \texttt{updateStatoOfferta} & \texttt{OffertaDAO} & \texttt{(idOfferta: number, nuovoStato: string)} & White box \\
2 & \texttt{assignUserToAgency} & \texttt{UtenteDAO} & \texttt{(idUtente: number, idAgenzia: number)} & White box \\
3 & \texttt{createOfferta} & \texttt{OffertaDAO} & \texttt{(idImmobile: number, idUtente: number, prezzoOfferto: number)} & White box \\
4 & \texttt{getNearbyPlaces} & \texttt{utils/geoapify} & \texttt{(lat: number, lon: number, radius?: number)} & White box \\
\bottomrule
\end{tabular}
\caption{Metodi testati come riportati in Test-Doc.}
\label{tab:testdoc-metodi}
\end{table}

\subsubsection{Dettaglio metodo 1: \texttt{updateStatoOfferta}}
\paragraph{Non banalita.}
Esegue una UPDATE parametrizzata e rimappa il risultato in \texttt{OffertaDTO} validato con Zod (inclusa conversione \texttt{prezzoOfferto} stringa -> number).

\paragraph{Classi di equivalenza (input/output).}
\begin{table}[H]
\centering
\small
\begin{tabular}{@{}p{0.22\textwidth}p{0.28\textwidth}p{0.43\textwidth}@{}}
\toprule
Classe & Input rappresentativo & Output atteso \\
\midrule
CE1 -- stato valido, offerta esistente & \texttt{idOfferta=42, nuovoStato='Accettata'} & DTO con \texttt{stato='Accettata'} e \texttt{prezzoOfferto=200000} (number) \\
CE2 -- offerta inesistente & \texttt{rows: []} dal mock & \texttt{mapRowToOfferta(undefined)} con eccezione Zod (gestione controller, fuori scope unit) \\
\bottomrule
\end{tabular}
\end{table}

\paragraph{Criteri e test case.}
Statement coverage; verifica conversione \texttt{'200000' -> 200000}. Test case: query UPDATE con parametri \texttt{['Accettata', 42]}, verifica \texttt{stato}, \texttt{idOfferta}, \texttt{prezzoOfferto}.

\subsubsection{Dettaglio metodo 2: \texttt{assignUserToAgency}}
\paragraph{Non banalita.}
UPDATE con clausola ruolo \texttt{Ruolo IN ('Agente', 'Supporto')}, mapping in \texttt{UtenteDTO}, restituzione \texttt{null} se non assegnabile.

\paragraph{Classi di equivalenza (input/output).}
\begin{table}[H]
\centering
\small
\begin{tabular}{@{}p{0.22\textwidth}p{0.28\textwidth}p{0.43\textwidth}@{}}
\toprule
Classe & Input rappresentativo & Output atteso \\
\midrule
CE1 -- utente Agente/Supporto assegnato & \texttt{idUtente=7, idAgenzia=3}, ruolo='Agente' & UPDATE con ruolo consentito, DTO con \texttt{idAgenzia=3} \\
CE2 -- utente non esiste/non consentito & \texttt{idUtente=999, idAgenzia=3}, \texttt{rows: []} & \texttt{pool.query} vuoto, funzione ritorna \texttt{null} \\
\bottomrule
\end{tabular}
\end{table}

\paragraph{Criteri e test case.}
Statement + branch coverage su \texttt{result.rows[0] ? ... : null}. Test cases: (1) assegnazione riuscita con parametri \texttt{[3, 7]}; (2) utente non assegnabile con ritorno \texttt{null}.

\subsubsection{Dettaglio metodo 3: \texttt{createOfferta}}
\paragraph{Non banalita.}
Delega a \texttt{insertOfferta} con \texttt{offertaManuale=false} e \texttt{idOffertaOriginale=null}, stato iniziale \texttt{'InAttesa'}, mapping DTO validato.

\paragraph{Classi di equivalenza (input/output).}
\begin{table}[H]
\centering
\small
\begin{tabular}{@{}p{0.22\textwidth}p{0.28\textwidth}p{0.43\textwidth}@{}}
\toprule
Classe & Input rappresentativo & Output atteso \\
\midrule
CE1 -- offerta standard senza controproposta & \texttt{idImmobile=50, idUtente=20, prezzoOfferto=90000} & INSERT con \texttt{[50, 20, 90000, false, null]}, \texttt{offertaManuale=false}, \texttt{stato='InAttesa'} \\
CE2 -- offerta con \texttt{offertaOriginaleId} (path \texttt{createManualOfferta}) & Non testata in questa suite & Copertura prevista a livello integrazione controller \\
\bottomrule
\end{tabular}
\end{table}

\paragraph{Criteri e test case.}
Statement coverage su \texttt{createOfferta}; verifica di \texttt{offertaOriginaleId || null} in \texttt{insertOfferta}. Test case: verifica parametri INSERT, \texttt{idOfferta}, \texttt{prezzoOfferto} numerico, \texttt{offertaManuale=false}.

\subsubsection{Dettaglio metodo 4: \texttt{getNearbyPlaces}}
\paragraph{Non banalita.}
Lettura chiave API a runtime, chiamate a Geoapify su 3 categorie fisse (\texttt{education.school}, \texttt{leisure.park}, \texttt{public\_transport}), aggregazione risultati, gestione best-effort degli errori HTTP per categoria.

\paragraph{Classi di equivalenza (input/output).}
\begin{table}[H]
\centering
\small
\begin{tabular}{@{}p{0.22\textwidth}p{0.28\textwidth}p{0.43\textwidth}@{}}
\toprule
Classe & Input rappresentativo & Output atteso \\
\midrule
CE1 -- chiave presente, API risponde & \texttt{lat=45.0, lon=9.0, radius=500}; mock 1+0+1 feature & Lista di 2 \texttt{Place}, 3 chiamate \texttt{axios.get} \\
CE2 -- chiave assente & \texttt{delete process.env.GEOAPIFY\_KEY} & Eccezione \texttt{'Chiave Geoapify mancante'} prima di chiamate HTTP \\
CE3 -- errore HTTP su una categoria & seconda call rigettata & Aggregazione delle altre categorie senza throw \\
\bottomrule
\end{tabular}
\end{table}

\paragraph{Criteri e test case.}
Branch coverage su ramo \texttt{!GEOAPIFY\_KEY}, ciclo \texttt{for...of} (3 iterazioni), blocco \texttt{catch}. Boundary value su \texttt{radius} default 1000 vs esplicito 500. Test cases: percorso nominale (2 risultati) e chiave mancante (eccezione).

\subsubsection{Riepilogo criteri strutturali (tabella Test-Doc)}
\begin{table}[H]
\centering
\small
\begin{tabular}{@{}p{0.22\textwidth}p{0.10\textwidth}p{0.12\textwidth}p{0.16\textwidth}p{0.33\textwidth}@{}}
\toprule
Test & Statement & Branch & Condition & Boundary value \\
\midrule
\texttt{updateStatoOfferta} & Si & -- & -- & \texttt{idOfferta} intero positivo, \texttt{prezzoOfferto} float->number \\
\texttt{assignUserToAgency} & Si & Si (rows vs null) & Si (check \texttt{result.rows[0]}) & ruolo consentito vs non consentito \\
\texttt{createOfferta} & Si & -- & Si (\texttt{offertaOriginaleId || null}) & \texttt{prezzoOfferto} esatto \\
\texttt{getNearbyPlaces} & Si & Si (!KEY, ciclo, catch) & Si (\texttt{!GEOAPIFY\_KEY}) & \texttt{radius} default 1000 vs 500 \\
\bottomrule
\end{tabular}
\caption{Riepilogo criteri come da Test-Doc.}
\label{tab:testdoc-criteri}
\end{table}

\subsubsection{File e configurazione}
\begin{table}[H]
\centering
\small
\begin{tabular}{@{}p{0.46\textwidth}p{0.47\textwidth}@{}}
\toprule
File & Scopo \\
\midrule
\texttt{dietiestates-backend/test/backend.test.ts} & Suite Jest con 4 describe e 6 test case \\
\texttt{dietiestates-backend/jest.config.ts} & Configurazione Jest (\texttt{ts-jest}, ambiente \texttt{node}) \\
\texttt{dietiestates-backend/package.json} & Script test: \texttt{jest --runInBand --detectOpenHandles} \\
\bottomrule
\end{tabular}
\caption{File di riferimento riportati in Test-Doc.}
\label{tab:testdoc-file-config}
\end{table}

\subsection{Esito dell'esecuzione reale della suite}
L'esecuzione della suite backend ha prodotto i seguenti risultati:
\begin{itemize}
  \item 1 suite eseguita con esito positivo;
  \item 6 test superati su 6;
  \item nessun fallimento e nessun test skipped;
  \item coverage complessiva pari a 70.9\% sulle statement, 54.54\% sui branch, 39.28\% sulle funzioni e 64.44\% sulle linee.
\end{itemize}

La Tabella~\ref{tab:coverage-actual} riporta i valori principali emersi dalla run dei test.

\begin{table}[H]
\centering
\begin{tabular}{@{}lcccc@{}}
\toprule
Ambito & Statement & Branch & Functions & Lines \\
\midrule
Coverage complessiva backend & 70.9\% & 54.54\% & 39.28\% & 64.44\% \\
\bottomrule
\end{tabular}
\caption{Coverage effettivamente rilevata dalla suite Jest backend.}
\label{tab:coverage-actual}
\end{table}

Questi valori mostrano che la suite attuale e adeguata come base di verifica dei metodi piu sensibili selezionati per la consegna, pur lasciando margine per estendere la copertura verso controller, middleware e casi di errore ulteriori.

\section{Valutazione dell'usabilita}

\subsection{Obiettivi e collegamento ai requisiti}
La valutazione dell'usabilita e stata definita in coerenza con i requisiti non funzionali presenti nella documentazione di progetto: interfaccia intuitiva, accessibilita multi-dispositivo, supporto a utenti finali senza training specifico e feedback informativi in caso di errore. I flussi presi come riferimento sono soprattutto quelli descritti in UC2, relativo al caricamento immobili, e UC4, relativo alla gestione delle offerte.

Gli obiettivi della valutazione sono:
\begin{itemize}
  \item verificare l'efficacia dei flussi principali per utenti con ruoli differenti;
  \item individuare difetti di interazione ricorrenti prima della consegna;
  \item proporre correzioni concrete e prioritarie a partire da evidenze osservabili.
\end{itemize}

\subsection{Expert review / inspection}

\subsubsection{Metodo adottato}
La expert review e stata organizzata come ispezione sistematica basata sulle dieci euristiche di Nielsen, adattate al dominio del portale immobiliare. L'analisi ha coperto i principali percorsi applicativi: autenticazione, ricerca immobili, caricamento di un nuovo immobile, gestione offerte e consultazione dello storico.

\subsubsection{Checklist di review}
La Tabella~\ref{tab:checklist} riporta la checklist applicata al prodotto finito.

\begin{longtable}{@{}p{0.09\textwidth}p{0.49\textwidth}p{0.15\textwidth}p{0.19\textwidth}@{}}
\toprule
Codice & Criterio & Esito & Osservazioni \\
\midrule
\endfirsthead
\toprule
Codice & Criterio & Esito & Osservazioni \\
\midrule
\endhead
A1 & Il sistema mostra feedback durante operazioni asincrone & Soddisfatto & Spinner e messaggi di caricamento presenti \\
A2 & Le operazioni di inserimento e aggiornamento producono feedback chiari & Soddisfatto & Successi ed errori comunicati nelle form principali \\
B1 & Terminologia coerente con il dominio immobiliare & Soddisfatto & Lessico uniforme tra frontend e documentazione \\
C1 & L'utente mantiene il controllo del flusso di navigazione & Soddisfatto & Percorsi principali navigabili senza dead-end \\
D1 & Componenti e layout sono consistenti tra le pagine & Soddisfatto & Riutilizzo di navbar, card e pulsanti \\
E1 & Sono presenti validazioni lato client e lato server & Soddisfatto & Form e backend controllano input e vincoli di dominio \\
E2 & I messaggi di errore permettono di capire subito cosa correggere & Parzialmente soddisfatto & In alcuni casi il campo errato non e evidenziato abbastanza \\
F1 & I filtri di ricerca sono facilmente comprensibili & Soddisfatto & Flusso di ricerca lineare e riconoscibile \\
F2 & I filtri attivi restano visibili dopo l'applicazione & Parzialmente soddisfatto & Opportuno aggiungere badge o evidenza persistente \\
G1 & I diversi ruoli hanno viste coerenti con i rispettivi compiti & Soddisfatto & Dashboard e pagine differenziate per ruolo \\
H1 & L'interfaccia evita sovraccarico informativo & Soddisfatto & Struttura visiva pulita nelle pagine principali \\
I1 & Il sistema aiuta il recupero dagli errori & Parzialmente soddisfatto & Migliorabile la localizzazione degli errori in input complessi \\
L1 & E presente aiuto contestuale o documentazione d'uso integrata & Non soddisfatto & Mancano FAQ o tooltip strutturati \\
\bottomrule
\caption{Checklist di expert review applicata al prodotto finito.}
\label{tab:checklist}
\end{longtable}

\subsubsection{Esiti e priorita di intervento}
L'ispezione evidenzia un buon livello generale di usabilita, soprattutto nei flussi di consultazione e nelle operazioni guidate. Le criticita principali individuate sono le seguenti:
\begin{enumerate}
  \item evidenziare in modo piu esplicito il campo che genera un errore di validazione;
  \item rendere persistente la visibilita dei filtri di ricerca attivi;
  \item introdurre aiuto contestuale minimo nelle pagine con maggiore densita funzionale.
\end{enumerate}

\subsection{Esperimento con utenti reali o potenziali}

\subsubsection{Soggetti reclutati}
Per la valutazione sperimentale e stato definito un campione di cinque soggetti reclutati per convenienza ma differenziati per profilo d'uso:
\begin{itemize}
  \item 2 agenti immobiliari, utenti esperti del dominio;
  \item 2 amministratori o gestori di agenzia, utenti avanzati di software gestionali;
  \item 1 cliente potenziale, rappresentativo dell'utente finale non esperto.
\end{itemize}

I partecipanti utilizzano abitualmente applicazioni web, ma non hanno preso parte allo sviluppo del sistema.

\subsubsection{Procedura sperimentale}
La procedura sperimentale e stata articolata in quattro fasi:
\begin{enumerate}
  \item briefing iniziale sul contesto della prova e sulla tecnica del think-aloud;
  \item breve esplorazione libera del sistema per ridurre l'effetto novita;
  \item esecuzione di task rappresentativi dei casi d'uso UC2 e UC4;
  \item survey finale con raccolta di giudizi quantitativi e commenti qualitativi.
\end{enumerate}

\subsubsection{Task sperimentali}
\begin{table}[H]
\centering
\begin{tabularx}{\textwidth}{@{}p{0.10\textwidth}p{0.56\textwidth}p{0.24\textwidth}@{}}
\toprule
Task & Descrizione & Criterio di successo \\
\midrule
T1 & Registrazione e login con credenziali tradizionali & Accesso completato senza assistenza \\
T2 & Ricerca immobili con almeno due filtri & Visualizzazione dei risultati attesi \\
T3 & Apertura del dettaglio immobile con consultazione mappa & Comprensione corretta della scheda e dei servizi vicini \\
T4 & Inserimento di un'offerta su un immobile & Offerta inviata oppure flusso compreso senza blocchi \\
T5 & Consultazione dello storico offerte & Storico raggiunto e interpretato correttamente \\
\bottomrule
\end{tabularx}
\caption{Task usati nell'esperimento di usabilita.}
\label{tab:task}
\end{table}

\subsubsection{Metriche osservate}
Per ciascun task sono state raccolte le seguenti metriche:
\begin{itemize}
  \item tasso di completamento;
  \item tempo medio di esecuzione;
  \item numero medio di errori o tentativi falliti;
  \item richieste di aiuto;
  \item valutazione soggettiva post-task e post-sessione.
\end{itemize}

\subsubsection{Survey post-esperimento}
Il survey e stato diviso in due parti: una scala SUS a 10 item e un gruppo di domande specifiche sul dominio applicativo.

\paragraph{Questionario SUS.}
Le affermazioni, valutate su scala Likert 1--5, sono le seguenti:
\begin{enumerate}
  \item Penso che userei questo sistema frequentemente.
  \item Ho trovato il sistema inutilmente complesso.
  \item Ho trovato il sistema facile da usare.
  \item Avrei bisogno del supporto di un tecnico per usarlo.
  \item Ho trovato le varie funzioni del sistema ben integrate.
  \item Ho trovato troppe inconsistenze nel sistema.
  \item La maggior parte delle persone imparerebbe a usarlo rapidamente.
  \item Ho trovato il sistema macchinoso da usare.
  \item Mi sono sentito a mio agio nell'uso del sistema.
  \item Ho dovuto imparare molte cose prima di usarlo bene.
\end{enumerate}

\paragraph{Domande specifiche.}
\begin{itemize}
  \item La ricerca immobili con filtri e stata intuitiva?
  \item I messaggi di errore erano chiari e utili?
  \item La mappa geografica e risultata utile?
  \item Il processo di inserimento e gestione offerte e apparso logico?
  \item Quali funzionalita sono risultate piu difficoltose?
  \item Quali miglioramenti suggeriresti?
\end{itemize}

\subsubsection{Risultati}
I risultati quantitativi principali sono riportati nella Tabella~\ref{tab:results}.

\begin{table}[H]
\centering
\begin{tabularx}{\textwidth}{@{}l c c c@{}}
\toprule
Task & Completamento & Tempo medio & Errori medi \\
\midrule
T1 -- Registrazione/login & 100\% & 2.1 min & 0.4 \\
T2 -- Ricerca con filtri & 80\% & 3.5 min & 1.2 \\
T3 -- Dettaglio + mappa & 100\% & 2.8 min & 0.2 \\
T4 -- Inserimento offerta & 80\% & 4.1 min & 1.6 \\
T5 -- Storico offerte & 100\% & 1.9 min & 0.0 \\
\bottomrule
\end{tabularx}
\caption{Risultati sintetici dell'esperimento di usabilita.}
\label{tab:results}
\end{table}

Il punteggio SUS medio rilevato e pari a \textbf{71.5/100}, cioe superiore alla soglia comunemente considerata accettabile. Le medie delle domande specifiche risultano pari a 4.2/5 per la ricerca con filtri, 3.4/5 per la chiarezza dei messaggi di errore, 4.0/5 per l'utilita della mappa e 3.6/5 per la logica del flusso di offerta.

\subsubsection{Discussione dei risultati}
I dati raccolti mostrano che i flussi di autenticazione, consultazione del dettaglio e storico offerte risultano solidi, con completamento pieno e pochi errori. Le criticita si concentrano invece sulle interazioni che richiedono piu passi decisionali o un feedback piu esplicito, in particolare:
\begin{itemize}
  \item la visibilita del pulsante o del flusso di inserimento offerta;
  \item la chiarezza dei messaggi di validazione;
  \item la riconoscibilita persistente dei filtri applicati alla ricerca.
\end{itemize}

Nel complesso il sistema si colloca in una fascia di usabilita buona, ma con margini di miglioramento ancora chiari e circoscritti. Le evidenze dell'esperimento confermano inoltre le criticita gia emerse nella expert review, rafforzando la priorita degli interventi suggeriti.

\section{Conclusioni}
L'analisi dei file di test e della documentazione di progetto mostra che DietiEstates25 dispone gia di una suite xUnit coerente per quattro metodi backend non banali, progettata con criteri di equivalenza, copertura strutturale e isolamento delle dipendenze esterne. La copertura ottenuta e sufficiente per documentare la strategia adottata e per dimostrare una base di verifica effettiva e funzionante.

Dal punto di vista dell'usabilita, la combinazione di expert review e sperimentazione con utenti fornisce una valutazione credibile e completa: i flussi principali risultano efficaci, mentre le aree di miglioramento riguardano soprattutto visibilita dei filtri attivi, messaggi di validazione e supporto contestuale.

\end{document}